%% discussion.tex

\section{Discussion: Architectural Design Recommendations}
\label{sec:discussion}

The evidence from our testbed validation (\S\ref{sec:evaluation}) and the cross-framework structural analysis (\S\ref{sec:related-work}) reveals a systematic architectural pattern: self-healing LLM agent systems consistently treat the recovery sink as a trusted channel, exempting it from the privilege invariants enforced on the forward execution path. In this section, we synthesize our findings into actionable design recommendations for framework developers, framed in terms of software architecture patterns and anti-patterns.

\subsection{Anti-Pattern: The Unchecked Recovery Sink}

The central anti-pattern we identify is the \emph{unchecked recovery sink}: a recovery pipeline that consumes a diagnostic LLM's recommendation (such as a role suggestion) and executes a privileged side effect (such as role replacement) without re-checking the privilege invariant that gates the forward execution path. Our cross-framework analysis (\S\ref{sec:related-work}) found this anti-pattern in production frameworks that support topology mutation---AutoGen~\cite{autogen} supports dynamic role reassignment, and \system{}~\cite{omo} ships a \code{TopologyMutationStrategy} that consults a diagnostic LLM. In each case, the recovery orchestrator exercises role-replacement authority on the basis of an LLM diagnosis whose input may embed attacker-influenced content.

The root cause of this anti-pattern is an \emph{implicit privilege assumption} inherited from classical OS exception handling: ``failure is benign, recovery action is safe.'' This assumption is valid when failure payloads and diagnostic outputs are generated by the trusted computing base (as in classical kernels), but breaks down when the diagnostic LLM's input (the semantic core dump) may contain external content that the agent was processing when it crashed. The fix is not to abandon self-healing---which provides well-documented robustness benefits---but to \emph{port the corresponding OS privilege invariant} to the recovery path.

\subsection{Design Pattern: The \reauthgate{}}

\paragraph{Intent.} Ensure that recovery actions never exceed the privilege of the failed execution they remedy.

\paragraph{Structure.} Centralize privilege re-verification at a single choke point (the \reauthgate{}) through which all side-effecting recovery actions must flow. Recovery strategies declare the need for reauthorization via a flag on the result type; the gate enforces $P(\text{recovery action}) \leq P(\text{failed execution})$ before any side effect executes. This is a \emph{prevent} control, not a \emph{detect-and-alert} control: unauthorized actions are blocked before state changes occur.

\paragraph{OS analogy.} This mirrors the OS principle that an exception handler executes with the privileges of the faulting context, never higher. The same principle underlies the confused-deputy fix~\cite{hardy1988confused}: the privileged deputy may not exercise authority not granted to the caller.

\paragraph{Why centralize?} Centralizing reauthorization at the orchestrator yields four software-engineering advantages (detailed in \S\ref{sec:why-centralize-pattern}): (i)~unified policy, (ii)~prevent rather than detect, (iii)~separation of concerns (strategy authors need not be security experts), and (iv)~open/closed for extension (new strategies inherit protection automatically). This is a classic cross-cutting concern pattern~\cite{aspectj}: permission enforcement is an aspect that every side-effecting recovery action needs, but is orthogonal to the strategy's core logic.

\paragraph{Validation.} On the testbed, this pattern eliminates the topology-mutation attack \emph{by construction}---the whitelist invariant is a programmatic check, not a probabilistic defense. This guarantee is LLM-independent: because the gate is a programmatic check on a whitelist, it is expected to hold across all models, regardless of instruction-following strength or adversary adaptiveness.

\subsection{Structural vs.\ Content-Level Enforcement: An Architecture Decision}

A key architectural decision emerges from the validation: \emph{structural defenses (privilege invariants enforced at a centralized choke point) are LLM-independent and therefore robust to model scaling}, whereas content-level defenses (which depend on the LLM obeying an untrusted-framing warning) may degrade on stronger instruction-following models. By construction, the \reauthgate{}'s effectiveness does not vary with the model's instruction-following strength, because the gate does not inspect the content of the diagnostic LLM's output---it only checks whether the extracted role string is in the pre-approved set and whether it is a privilege escalation from the failing role.

This asymmetry has an architectural implication: \emph{as LLMs become stronger instruction-followers, structural enforcement transitions from complementary to necessary.} Framework designers should not rely on content-level provenance signaling alone (template selection, untrusted framing) but should pair it with structural privilege invariants at the recovery sink. Content-level defenses remain valuable as defense-in-depth---they may catch attacks that the structural invariant does not cover (e.g., content-level trust-escalation via high-trust template re-injection, which is outside the scope of this paper but is addressed in the broader agent-security literature~\cite{debenedetti2025camel,zhan2025adaptive})---but they should not be the sole layer of defense for privileged recovery actions.

\subsection{Recommendations for Framework Developers}

Based on our analysis, we recommend that developers of self-healing LLM agent frameworks:

\begin{enumerate}[leftmargin=*]
    \item \textbf{Audit recovery sinks for missing authorization.} Check whether the role-replacement / topology-mutation primitive performs an authorization check on the value it consumes, or whether it consumes the diagnostic LLM's recommendation verbatim. The CWE-862 criterion from \S\ref{sec:attack-chain} provides a replicable methodology: if the sink performs a privileged operation without verifying that the caller is authorized to request that operation, the sink is vulnerable.

    \item \textbf{Centralize privilege re-verification.} Do not scatter permission checks across individual recovery strategies. Centralize them at the orchestrator level, where the invariant ``no recovery action exceeds the privilege of the failed execution'' can be verified by inspecting one method. This locality is itself an OS principle: classical kernels enforce privilege invariants at well-defined boundaries (syscall entry, exception dispatch), not inside each driver or subsystem.

    \item \textbf{Defer side effects, do not execute inline.} Side-effecting recovery strategies should not execute their side effects inline at strategy-return time. Instead, they should package the side effect as a deferred callback, set a \code{requiresReauthorization} flag on the result, and let the orchestrator execute the callback only after the gate has verified the privilege invariant. This turns the security property from an emergent property of independently-authored strategies into a local, inspectable property of one orchestrator method.

    \item \textbf{Pair content-level with structural defenses.} Do not rely solely on prompt-level untrusted framing or instruction-hierarchy enforcement for privileged recovery actions. As models become stronger instruction-followers, content-level defenses degrade; structural privilege invariants provide a model-independent backstop.
\end{enumerate}

\subsection{Migration Guide for Framework Developers}
\label{sec:migration-guide}

We now provide a step-by-step migration guide for porting the \reauthgate{} design pattern onto an existing self-healing pipeline. The guide assumes a typical framework architecture with (i)~a tool-execution layer that raises exceptions, (ii)~a recovery orchestrator that catches failures and dispatches recovery strategies, and (iii)~a permission checker that gates forward-path tool calls. The migration is designed to be incremental and backward-compatible.

\paragraph{Step 1: Add a \code{requiresReauthorization} flag to the recovery result type.}
The first and lowest-risk step is to extend the recovery result type with a boolean flag \code{requiresReauthorization} (default \code{false}). No gate is added yet; this step only prepares the plumbing. Existing strategies that do not set the flag default to \code{false} and behave as before (no gate interception), so a framework can migrate strategies incrementally without breaking existing functionality.

\paragraph{Step 2: Defer side effects in side-effecting strategies.}
For each side-effecting recovery strategy (role replacement, tool execution, model switching), refactor the strategy to package its side effect as a \emph{deferred callback} stored in the recovery context metadata, rather than executing it inline at strategy-return time. The strategy sets \code{requiresReauthorization=true} and returns; the side effect is held in suspension. This step does not change observable behavior yet (no gate intercepts the flag), but it prepares the strategy for gate interception.

\paragraph{Step 3: Install the \reauthgate{} at the orchestrator.}
Modify the recovery orchestrator's main loop to intercept every result with \code{requiresReauthorization=true} \emph{before} executing the deferred callback. The gate calls the existing \code{PermissionChecker} (the same one used on the forward path) to verify the privilege invariant: for role replacement, the new role must be in the pre-approved downgrade set of the original role; for tool execution, the requested tool must be in the original role's allowed set. If the check fails, the side effect is blocked and the orchestrator returns a failed recovery result; if it passes, the deferred callback executes. The gate is a single method that can be inspected to verify the invariant ``no recovery action exceeds the privilege of the failed execution.''

\paragraph{Step 4: Add a role-existence whitelist.}
In addition to the privilege check, add a whitelist check: the proposed role must exist in the framework's role blueprint registry, not be an arbitrary LLM-suggested string. This closes the CWE-862 sink completely: even if the diagnostic LLM is tricked into suggesting an out-of-set role (e.g., \code{admin}), the whitelist rejects it before the side effect executes.

\paragraph{Step 5: Audit existing strategies for compliance.}
After installing the gate, audit each existing recovery strategy to ensure that (a)~every side-effecting strategy sets \code{requiresReauthorization=true}, and (b)~no strategy executes side effects inline (bypassing the gate). This audit is local: it checks one flag per strategy, not the strategy's internal logic. Strategies that are pure diagnostic (no side effects) need no change. New strategies added after the migration inherit protection automatically by setting the flag.

\paragraph{Expected migration cost.}
In the \system{} testbed, the full migration touched 14 strategy classes and one orchestrator method, with sub-millisecond runtime overhead and zero false alarms on benign recovery (\S\ref{sec:evaluation}). The migration is designed to be backward-compatible: a strategy that does not set \code{requiresReauthorization} defaults to \code{false} and behaves as before (no gate interception), so a framework can migrate strategies incrementally without breaking existing functionality. The pattern does not require any modification to the forward execution path, the LLM, or the role/permission model.
